\chapter{Supporting Material}
\label{appx:material}

This appendix reproduces the full captures that back the diagrams, tables and claims of the Design and Implementation chapter (Chapter~\ref{chap:execution}) and the Critical Evaluation chapter (Chapter~\ref{chap:evaluation}). The configuration and pipeline captures of Sections~\ref{appx:flows} to~\ref{appx:runtime} are from the stable baseline of 3 August 2026; the evaluation artefacts of Sections~\ref{appx:matrix} onward are from the dedicated fresh runs of 5 to 7 August 2026 that produced Chapter~\ref{chap:evaluation}. All were taken from the running testbed with the defence armed except where a disarmed contrast is stated. The short outputs appear inline in the chapters; the long ones are collected here in full.

\paragraph{Code and data availability.} The complete system, every evaluation harness and the raw result files are published at \href{\repourl}{the project GitHub repository}. The excerpts reproduced in this appendix are the load-bearing lines of larger files; the full versions, and the captured data behind every chart and table, are in the repository at the paths given in Table~\ref{tab:repomap}. Each path below is a live link to its file in the repository.

\begin{table}[h]
\centering
\small
\begin{tabular}{p{0.34\textwidth}p{0.58\textwidth}}
\hline
Repository path & Role in this dissertation \\
\hline
\href{\repourl/blob/main/06_Build/cars_engine.py}{\texttt{06\_Build/cars\_engine.py}} & The controller: \texttt{classify()}, criticality elevation, \texttt{select\_response()}, GUARD binding, reactive install (cookie \texttt{0x00ca}, \texttt{30+15w}), authenticated API \\
\href{\repourl/blob/main/06_Build/snort_bridge.py}{\texttt{06\_Build/snort\_bridge.py}} & Snort alert to \texttt{/cars/respond} detection bridge \\
\href{\repourl/blob/main/06_Build/cars_flow_audit.py}{\texttt{06\_Build/cars\_flow\_audit.py}} & Flow-integrity poll against the trusted baseline \\
\href{\repourl/blob/main/06_Build/cars_remediation.py}{\texttt{06\_Build/cars\_remediation.py}} & Process last-good restore agent \\
\href{\repourl/blob/main/06_Build/cars_process.py}{\texttt{06\_Build/cars\_process.py}} & Process / control-law reference model \\
\href{\repourl/tree/main/07_Evaluation/overnight}{\texttt{07\_Evaluation/overnight/}} & Conformance-at-scale, criticality behaviour, flow-integrity, jitter and the Gap 1/2/4 harnesses, with raw data under \texttt{results/} \\
\href{\repourl/blob/main/07_Evaluation/overnight/gap4_flowmonitor/cars_flowmonitor.py}{\texttt{gap4\_flowmonitor/cars\_flowmonitor.py}} & Event-driven flow-integrity monitor (Section~\ref{sec:threats}) \\
\href{\repourl/tree/main/report}{\texttt{report/}} & This dissertation, its figures and its \texttt{appendix\_listings/} data files \\
\hline
\end{tabular}
\caption{Map from the report to the published repository (\href{\repourl}{the project GitHub repository}).}
\label{tab:repomap}
\end{table}

\lstdefinestyle{carsdump}{basicstyle=\ttfamily\tiny,breaklines=true,breakatwhitespace=false,columns=fullflexible,frame=single,captionpos=b,showstringspaces=false,keepspaces=true,xleftmargin=3pt,xrightmargin=3pt}
\lstdefinestyle{carssrc}{basicstyle=\ttfamily\scriptsize,breaklines=true,breakatwhitespace=false,columns=fullflexible,frame=single,captionpos=b,showstringspaces=false,keepspaces=true}

\section{Full OpenFlow tables}
\label{appx:flows}

The complete three-table pipeline of Section~\ref{sec:pipeline}, dumped from each switch with \texttt{ovs-ofctl -O OpenFlow13 dump-flows}. Table~0 is GUARD (identity binding and anti-spoof), Table~1 is POLICY (the stateful conntrack layer with the allowlist under cookie \texttt{0xa2} and the default-deny), and Table~2 is the L2 learning switch. These are the source for the pipeline in Figure~\ref{fig:carsflow}.

\lstinputlisting[style=carsdump,caption={Cell-1 switch \texttt{ovs1}, tables 0 to 2.},label={lst:ovs1}]{appendix_listings/flows_ovs1.txt}
\lstinputlisting[style=carsdump,caption={Gateway switch \texttt{ovsgw}, tables 0 to 2.},label={lst:ovsgw}]{appendix_listings/flows_ovsgw.txt}
\lstinputlisting[style=carsdump,caption={Cell-2 switch \texttt{ovs2}, tables 0 to 2.},label={lst:ovs2}]{appendix_listings/flows_ovs2.txt}

\section{Deployed decision logic}
\label{appx:logic}

The rulebook as it is defined in the running \href{\repourl/blob/main/06_Build/cars_engine.py}{\texttt{cars\_engine.py}}, backing Table~\ref{tab:rulebook}. The ordering is significant: the engineering and remediation rows precede the dangerous-operations block, so their pre-authorised control of the PLC classifies as OPERATIONAL. This is the current deployed rulebook, in which \texttt{ews->plc} is OPERATIONAL.

\lstinputlisting[style=carssrc,caption={The deployed \texttt{RULEBOOK} (first match wins).},label={lst:rulebook}]{appendix_listings/rulebook_src.txt}

The proactive A2 policy read from \texttt{a2\_policy.json} on the controller: the nine allowlisted conduits and the default-deny set that back the allowlist and default-deny rows of Table~1.

\lstinputlisting[style=carssrc,caption={The A2 allowlist and default-deny (\texttt{a2\_policy.json}).},label={lst:a2}]{appendix_listings/a2_policy.txt}

\section{Detection rules}
\label{appx:snort}

The Snort rules that recover the industrial operation by deep packet inspection, backing Figures~\ref{fig:detresp} and~\ref{fig:carsflow}. Each rule maps an S7 or Modbus function to one of the operation classes CARS reasons about. The S7 rules are shown in their hardened form: after the fragmentation finding of Section~\ref{sec:evaldpi}, the operation match was re-anchored to the S7 job header with \texttt{distance}/\texttt{within} rather than a fixed packet offset, and TCP stream reassembly was enabled on the ICS ports (Listing~\ref{lst:dpihard}), so a PDU split across TCP segments is still recovered.

\lstinputlisting[style=carssrc,caption={S7 and Modbus operation-class rules from \texttt{cars.rules}.},label={lst:snort}]{appendix_listings/snort_rules.txt}

\subsection{Sensor coverage boundary}
\label{appx:sensorcoverage}

The operation-aware tier depends on what the Snort mirror carries, and that mirror taps the gateway switch alone. Figure~\ref{fig:cellblindspot} illustrates the spatial consequence noted in Section~\ref{sec:threats}: traffic that crosses \texttt{ovsgw} is inspected, but an attack that stays within a single cell, such as one from a foothold inside Cell-2 against \texttt{PLC2} on \texttt{ovs2}, never reaches the mirror and draws no operation-aware response. The proactive identity binding and default-deny run on every switch independently of the sensor, so an unregistered intra-cell source is still refused; the residual is the operation-aware tier for a compromised, already-allowlisted conduit. No intra-cell attacker was present in the evaluated set, so this is an architectural illustration rather than a measured attack, and per-cell sensing is the mitigation carried in the future work.

\begin{figure}[t]
\centering
\includegraphics[width=\textwidth]{fig_cellblindspot.png}
\caption{The spatial coverage of the operation-aware tier. The Snort mirror taps \texttt{ovsgw} only: gateway-crossing traffic is inspected (green), while a purely intra-cell path (red), an in-cell foothold reaching \texttt{PLC2} over \texttt{ovs2}, never reaches the mirror and is not operation-aware. The proactive layer still applies on every switch, so an unregistered intra-cell source is refused regardless. This is an architectural illustration of a coverage boundary, not a tested attack.}
\label{fig:cellblindspot}
\end{figure}

\section{Runtime state}
\label{appx:runtime}

Two further live checks referenced from Section~\ref{sec:implementation}. Figure~\ref{fig:appguard} is the GUARD anti-spoof drop counter for every binding on every switch, all at zero under normal operation. Figure~\ref{fig:appmeter} is the rate-limiting meter that backs the THROTTLE response.

\begin{figure}[h]
\centering
\includegraphics[width=\textwidth]{term_guard.png}
\caption{GUARD anti-spoof drop counters from \texttt{curl \$API/cars/guard}, per binding and per switch (all zero, armed and quiescent).}
\label{fig:appguard}
\end{figure}

\begin{figure}[h]
\centering
\includegraphics[width=\textwidth]{term_meter.png}
\caption{The OpenFlow meter behind the THROTTLE response, from \texttt{ovs-ofctl dump-meters ovsgw}.}
\label{fig:appmeter}
\end{figure}

\section{Complete evaluation matrix}
\label{appx:matrix}

The full decision matrix of Section~\ref{sec:evalaccuracy}, all twenty-seven cases with the columns the body table abbreviates (source, destination, operation, rate, ground-truth label, tier, response, verdict, note), regenerated from the deployed engine via \texttt{/cars/respond}.

\lstinputlisting[style=carssrc,caption={The evaluation matrix, \texttt{cars\_eval\_matrix.csv}.},label={lst:matrix}]{appendix_listings/eval_matrix.csv}

\section{Decision logs and events}
\label{appx:logs}

The controller logs every decision it takes, so the conformance result rests on thousands of live judgements, not only the scored matrix. Listing~\ref{lst:decisions} is a representative extract spanning the campaign: the reconnaissance isolate, the identity-spoof refusal, the mean-time-to-mitigate isolate, the operation-aware isolate, the flow-integrity drift, and two legitimate operations permitted. The complete per-run logs are the exported CSV artefacts (each run logs on the order of a thousand decisions).

\lstinputlisting[style=carssrc,caption={Representative controller decisions across the campaign (extract of the decision log; times \texttt{08-05}).},label={lst:decisions}]{appendix_listings/decision_extract.csv}

\section{Packet-level and wire-level analysis}
\label{appx:wire}

The wire evidence is dissected directly from the captured pcaps. On the S7 side, the attacker's Write-Var frames to PLC1 were counted at the two synchronised capture points, the Snort mirror and the PLC's physical port, armed and disarmed, giving the drop-on-the-wire result of Table~\ref{tab:wiredrop}: eight write attempts seen by the detector, one reaching the PLC when armed against seven when disarmed.

On the Modbus side, the operation-awareness discriminator is visible in the frame bytes. The same conduit \texttt{192.168.2.31 $\to$ 192.168.2.20:502} carried a read and a write, distinguished only by the function code at offset seven of the payload:

\begin{lstlisting}[style=carssrc,caption={Modbus read and write on the same conduit; the function code sits at offset 7.},label={lst:modbusfc}]
READ  (-> ALLOW)     MBAP+PDU: 00 01 00 00 00 06 00 03 00 00 00 04
                     tid=0001 proto=0000 len=0006 unit=00 FC=03  (read holding)
WRITE (-> THROTTLE)  MBAP+PDU: 00 01 00 00 00 06 00 06 00 04 00 07
                     tid=0001 proto=0000 len=0006 unit=00 FC=06  (write register)
\end{lstlisting}

\noindent CARS allowed the FC\,03 read and throttled the FC\,06 write on that single byte. The dangerous function codes (\texttt{0x05} control, \texttt{0x08} diagnostic, \texttt{0x2B} program, \texttt{0x64} illegal) and the reliability with which the sensor detects each are given in Table~\ref{tab:dpi}. Figure~\ref{fig:modbusdpi} shows the two Modbus function-code fields as Wireshark dissects them.

\begin{figure}[h]
\centering
\begin{minipage}[t]{0.49\textwidth}\centering\includegraphics[width=\linewidth]{wire_modbus_fc3.png}\end{minipage}\hfill
\begin{minipage}[t]{0.49\textwidth}\centering\includegraphics[width=\linewidth]{wire_modbus_fc6.png}\end{minipage}
\caption{The Modbus operation on the wire, same conduit \texttt{.31$\to$.20}. Left: \texttt{Function Code: Read Holding Registers (3)}, permitted. Right: \texttt{Function Code: Write Single Register (6)} setting register~4 to~7, throttled. The verdict turns on this field.}
\label{fig:modbusdpi}
\end{figure}

Two further S7 views complete the picture. Figure~\ref{fig:wiredisarmed} shows the attacker's write reaching the PLC with CARS disarmed (contrast the single armed write of Figure~\ref{fig:wireproof}). Figure~\ref{fig:s7stream} shows that lone attacker frame amongst the legitimate Factory IO loop traffic on the PLC port, the attack is a needle in the live process stream.

\begin{figure}[h]
\centering
\includegraphics[width=0.9\textwidth]{wire_disarmed.png}
\caption{Disarmed, the attacker's S7 Write Var (frame 452) reaches the PLC. With CARS armed the same write is dropped in the fabric (Figure~\ref{fig:wireproof}).}
\label{fig:wiredisarmed}
\end{figure}

\begin{figure}[h]
\centering
\includegraphics[width=\textwidth]{wire_s7_stream.png}
\caption{The attacker's Setup frame (\texttt{.31$\to$.10}, highlighted) amongst the legitimate \texttt{.55} Factory IO loop traffic on the PLC port. CARS must pick the malicious operation out of a continuous stream of legitimate S7 reads and writes.}
\label{fig:s7stream}
\end{figure}

The enforcement itself is visible on a third wire, the OpenFlow control channel (TCP~6653), captured in \texttt{of\_control.pcap}. Listing~\ref{lst:flowmod} is the \texttt{OFPT\_FLOW\_MOD} that installs the reactive isolate, decoded from that capture: an ADD of the drop rule stamped with the CARS reactive cookie \texttt{0x00ca}, at priority~110 in the POLICY table, matching the attacker's source with a 75-second self-healing timeout. Two identical messages are sent, one to \texttt{ovs1} and one to \texttt{ovsgw}, so the quarantine is installed fabric-wide from a single decision. This is the control-plane counterpart to the datapath dump of Figure~\ref{fig:p6armedflow}: the rule seen being installed, not just seen in place.

\begin{lstlisting}[style=carssrc,caption={The reactive isolate as an \texttt{OFPT\_FLOW\_MOD}, decoded from \texttt{of\_control.pcap} (OpenFlow 1.3, tcp/6653).},label={lst:flowmod}]
OFPT_FLOW_MOD              (version 0x04, type 14, control channel tcp/6653)
  command      = ADD
  cookie       = 0x00000000000000ca     (CARS reactive marker; 0xa2 = static A2 policy)
  table_id     = 1                        (POLICY)
  priority     = 110                       (above the allowlist and default-deny)
  hard_timeout = 75 s                      (CRITICAL asset: 30 + 15*3, then self-heals)
  match        = eth_type=0x0800, ipv4_src=192.168.2.31
  instructions = APPLY_ACTIONS []          (no output action = drop)
-> sent to ovs1 and ovsgw from one controller decision (fabric-wide quarantine)
\end{lstlisting}

\section{Scenario testing in full}
\label{appx:scenarios}

The process devastation of Section~\ref{sec:evaldevastation} has a second, quieter variant: instead of driving the tank visibly over the top, the attacker pins the reported level near zero, so the operator panel shows a safe, empty tank while the real one fills. Figure~\ref{fig:p6deception} captures this deception from a disarmed run: the real Factory IO tank stands filled while the HMI reads a level of zero and a pump of \texttt{+0.000}.

\begin{figure}[h]
\centering
\begin{minipage}[c]{0.42\textwidth}\centering\includegraphics[width=\linewidth]{p6_tank_disarmed.png}\end{minipage}\hfill
\begin{minipage}[c]{0.54\textwidth}\centering\includegraphics[width=\linewidth]{p6_hmi_deceived.jpg}\end{minipage}
\caption{The sensor-spoof deception, disarmed. Left: the real Factory IO tank, filled. Right: the operator HMI at the same moment, reading zero. The reported process has been divorced from the real one.}
\label{fig:p6deception}
\end{figure}

The external vectors were driven from the Kali attacker \texttt{192.168.2.77}. Reconnaissance under CARS is shown in Figure~\ref{fig:p1recon}: an armed port scan of the PLC returns every port \texttt{filtered}, and the scan is reactively isolated with a criticality-scaled quarantine. Identity spoofing is shown in Figure~\ref{fig:p2spoof}: the GUARD anti-spoof drop counters register the forged ARP frames dropped at ingress. The reaction window over one hundred autonomous trials is in Figure~\ref{fig:mttm}.

\begin{figure}[h]
\centering
\includegraphics[width=0.82\textwidth]{p1_recon_ports.png}\\[4pt]
\includegraphics[width=0.82\textwidth]{p1_recon_isolate.png}
\caption{Reconnaissance from Kali, armed. Top: the port scan of PLC1, every port \texttt{filtered}. Bottom: the reactive rule installed against the scanner, \texttt{cookie=0xca}, \texttt{priority=110}, \texttt{nw\_src=192.168.2.77}, \texttt{actions=drop}, \texttt{hard\_timeout=75}.}
\label{fig:p1recon}
\end{figure}

\begin{figure}[h]
\centering
\includegraphics[width=\textwidth]{p2_guard_after.png}
\caption{Identity spoofing at GUARD. After forging ARP as the trusted seam \texttt{192.168.2.31} from Kali, the anti-spoof counter \texttt{dpid3 arp 192.168.2.31} registers the dropped frames (running total 10), every forged frame dropped at ingress with all other bindings untouched. The controller logged each as \texttt{REFUSED (identity-spoof)}.}
\label{fig:p2spoof}
\end{figure}

\begin{figure}[h]
\centering
\includegraphics[width=0.9\textwidth]{p2_mttm.png}
\caption{Reaction window (mean-time-to-mitigate) over one hundred autonomous trials, attack-frame to enforcement on a single clock: median 7.6\,ms, mean 11.2\,ms, standard deviation 11.8\,ms, minimum 5.5\,ms, maximum 72.9\,ms, every trial ending in ISOLATE.}
\label{fig:mttm}
\end{figure}

\section{Scripts and code}
\label{appx:code}

Each script is presented as its role plus its load-bearing snippet. The lead is the controller, \href{\repourl/blob/main/06_Build/cars_engine.py}{\texttt{cars\_engine.py}}: the whole evaluation drives its decision path, whose core is \texttt{classify()}, mapping a (source, destination, operation) to a tier by walking the rulebook first-match-wins.

\begin{lstlisting}[style=carssrc,caption={The controller decision path (\href{\repourl/blob/main/06_Build/cars_engine.py}{\texttt{cars\_engine.py}}): first-match rulebook classification.},label={lst:classify}]
def classify(src_ip, dst_ip, op=None):
    s, d = role_of(src_ip), role_of(dst_ip)      # map IP -> role from the REGISTRY
    for rsrc, rdst, rop, tier in RULEBOOK:        # rules in priority order
        if (rsrc == "any" or rsrc == s or rsrc == src_ip) and \
           (rdst == "any" or rdst == d or rdst == dst_ip) and \
           (rop  == "any" or rop  == op):
            return tier, s, d                     # first match wins -> its tier
    return "FORBIDDEN", s, d                       # nothing matched -> default deny
\end{lstlisting}

\noindent The role lookup resolves each address to its registered role; the loop returns the tier of the first rulebook row whose source, destination and operation all match, and defaults to FORBIDDEN when none do. This is the exact path the conformance evaluation drives through \texttt{/cars/respond}. The remaining load-bearing pieces of the controller and of the supporting scripts follow in the same form.

\subsection*{Controller: the rest of the decision-to-enforcement path}

\noindent \textbf{Criticality elevation and block duration (\texttt{respond()}).} After classification, the destination's criticality weight drives two things: a trusted actuating write (SENSITIVE) to a CRITICAL asset is elevated to FORBIDDEN unless an authorised maintenance window is open, and the quarantine duration is scaled by the same weight.

\begin{lstlisting}[style=carssrc,caption={Grey-zone elevation and criticality-scaled block duration.},label={lst:elevate}]
acl = crit_of(dst_ip); dcw = CW.get(acl, 0)          # dst asset weight: 3..0
in_window = time.time() < self.maint_until
elevated = (tier == "SENSITIVE" and dcw >= 3 and not in_window)
if elevated:
    tier = "FORBIDDEN"          # trusted WRITE to a CRITICAL asset -> forbidden
timeout = BLOCK_TIMEOUT + dcw * 15   # CRITICAL 75, HIGH 60, MED 45, LOW 30 s
\end{lstlisting}

\noindent \textbf{The graded response ladder (\texttt{select\_response()}).} A tier, the asset weight and a rate flag select one of the seven responses; the safety loop is never enforced against, and higher-criticality assets escalate sooner and skip rungs.

\begin{lstlisting}[style=carssrc,caption={Criticality-graded, flood-aware response selection.},label={lst:selresp}]
if tier == "CRITICAL":    return "REFUSE"        # safety invariant - never enforce
esc = max(1, ESCALATE - dcw); c = state.get("count", 1)
if tier == "OPERATIONAL":
    if flood:
        if dcw >= 3: return "BLOCK"              # flood on CRITICAL asset -> cut
        return "THROTTLE" if c < esc else "BLOCK"
    return "ALLOW"                                # trusted conduit, normal rate
if tier == "SENSITIVE":
    if flood:    return "BLOCK"
    if dcw >= 3: return "BLOCK"                   # CRITICAL asset: skip THROTTLE
    return "THROTTLE" if c < esc else "BLOCK"
if flood:    return "ISOLATE"                     # FORBIDDEN
if dcw >= 3: return "ISOLATE"                     # CRITICAL asset: isolate now
return "BLOCK" if src_count < esc else "ISOLATE"
\end{lstlisting}

\noindent \textbf{The reactive install (\texttt{isolate\_source()}).} Every enforcement is a real OpenFlow rule stamped with the distinct reactive cookie \texttt{0x00ca}, self-expiring via \texttt{hard\_timeout} and notifying the controller on removal (\texttt{SEND\_FLOW\_REM}) so the defence self-heals when the attack stops. Isolation drops all traffic from the source across every switch.

\begin{lstlisting}[style=carssrc,caption={Source quarantine: a self-healing drop on every switch (\texttt{REACTIVE\_COOKIE = 0x00CA}).},label={lst:isolate}]
def isolate_source(self, src_ip, timeout=BLOCK_TIMEOUT):
    for d, dp in list(self.datapaths.items()):
        ofp, psr = dp.ofproto, dp.ofproto_parser
        m = psr.OFPMatch(eth_type=0x0800, ipv4_src=src_ip)      # any dst
        dp.send_msg(psr.OFPFlowMod(datapath=dp, cookie=REACTIVE_COOKIE,
            table_id=1, priority=110, match=m, instructions=[],  # [] = drop
            hard_timeout=timeout, flags=ofp.OFPFF_SEND_FLOW_REM))
\end{lstlisting}

\noindent \textbf{Anti-spoof GUARD (\texttt{install\_guard()}).} Table~0 binds each protected identity to its physical port and MAC; a frame carrying a protected source address that does not arrive on its bound port is dropped before it can reach policy. This is prevention, not detection.

\begin{lstlisting}[style=carssrc,caption={Table 0 identity binding and anti-spoof drop.},label={lst:guard}]
for (d, port, mac, ip) in BINDINGS:                # bind identity to port+MAC
    if d == dp.id:
        add(200, OFPMatch(in_port=port, eth_type=0x0800,
                          eth_src=mac, ipv4_src=ip), GOTO)   # legit -> policy
if GUARD_ENABLED:
    for ip in PROTECTED_IPS:
        add(100, OFPMatch(eth_type=0x0800, ipv4_src=ip), [])  # forged src -> drop
add(0, OFPMatch(), GOTO)                            # default -> policy table
\end{lstlisting}

\noindent \textbf{Control-API authentication (\texttt{\_app()}).} Every state-changing endpoint requires a shared token; the read-only status endpoints and the Snort detection feed are excluded so the bridge is unchanged.

\begin{lstlisting}[style=carssrc,caption={Token gate on the control endpoints.},label={lst:auth}]
_CONTROL = ('/cars/defense','/cars/maintenance','/cars/reload','/cars/reload-a2',
            '/cars/block','/cars/unblock','/cars/restore')
if path in _CONTROL and environ.get('HTTP_X_CARS_TOKEN','') != self.api_token:
    return reply({"error": "unauthorized - X-CARS-Token required"}, '401 Unauthorized')
# /cars/respond (Snort bridge feed) is intentionally excluded
\end{lstlisting}

\subsection*{The supporting scripts}

\noindent \textbf{Detection bridge (\href{\repourl/blob/main/06_Build/snort_bridge.py}{\texttt{snort\_bridge.py}}).} Snort raises the alerts; the bridge recovers the ICS operation from the alert text, measures its arrival rate, and posts the source, destination, operation and rate to the controller, which alone decides. The brain, not the sensor, holds the policy.

\begin{lstlisting}[style=carssrc,caption={Operation recovery from the Snort alert, forwarded to the controller.},label={lst:bridge}]
API = "http://10.10.10.1:8080/cars/respond"     # the brain decides; bridge only reports
mo = re.search(r'CARS-(MODBUS|S7)-(WRITE|READ|CONTROL|DIAG|PROGRAM|ILLEGAL)', line)
if mo: proto = mo.group(1); op = mo.group(2)     # recover the operation
key = (src, dst, op)                             # keep read vs write distinct
# rate of this (src,dst,op) measured, then POSTed with op+rate to API
\end{lstlisting}

\noindent \textbf{Flow-integrity self-check (\href{\repourl/blob/main/06_Build/cars_flow_audit.py}{\texttt{cars\_flow\_audit.py}}).} A trusted baseline of the immutable policy (Table~0 GUARD and the Table~1 A2 allowlist, cookie \texttt{0xa2}) is captured at a clean start; the only additions permitted at runtime are the reactive rules (cookie \texttt{0xca}). Anything else appearing on a switch is flagged as drift and surfaced into the controller's decision log.

\begin{lstlisting}[style=carssrc,caption={Flow identity and the drift model.},label={lst:flowaudit}]
# identity = bridge,table,priority,cookie,MATCH (not actions): a tampered action on a
# policy rule shows as "changed"; a bogus new rule as "extra".
def key_of(br, s): return "%s|t%s|p%s|c%s|%s" % (br, s[0], s[1], s[2], s[3])
# baseline = GUARD (t0) + A2 allowlist (t1, cookie 0xa2); the ONLY legitimate live
# additions carry the reactive cookie 0xca. Anything else later -> drift -> audit.
\end{lstlisting}

\noindent \textbf{Process-state remediation (\href{\repourl/blob/main/06_Build/cars_remediation.py}{\texttt{cars\_remediation.py}}).} The novel maintain half of block-and-maintain: it detects sensor tampering by process anomaly (a level the bang-bang control law physically cannot produce) and writes the last-good value back over S7, holding the process while the source is quarantined.

\begin{lstlisting}[style=carssrc,caption={Anomaly-triggered last-good restore over S7.},label={lst:remediation}]
last_good = SEED
while True:
    lvl = s7_read(c)                         # read Tank.Level over S7
    if tampered(lvl, prev):                  # a value the control law cannot produce
        s7_write(c, last_good)               # RESTORE last-good on tamper
        feed("RESTORED", level=round(lvl,1), last_good=round(last_good,1))
    else:
        last_good = lvl                       # healthy -> advance last-good
\end{lstlisting}

\noindent \textbf{Latency harness (\href{\repourl/blob/main/07_Evaluation/cars2_mttm.py}{\texttt{cars2\_mttm.py}}).} Times the reaction window on a single clock: start the attack, poll for the isolate to appear in the controller's block set, and record attack-frame-to-enforcement.

\begin{lstlisting}[style=carssrc,caption={Single-clock mean-time-to-mitigate measurement.},label={lst:mttm}]
t0 = time.time()
# ... launch the attack ...
while not blocked() and time.time() - t0 < 15: time.sleep(0.02)  # poll for isolate
mttm = time.time() - t0                       # attack-frame -> enforcement
\end{lstlisting}

\noindent \textbf{Attack tool (\href{\repourl/blob/main/06_Build/cars_fdi_overflow.py}{\texttt{cars\_fdi\_overflow.py}}).} The false-data injection used in Section~\ref{sec:evalcampaign}: a compromised host opens an S7 session to the PLC and writes the output area at high rate, blinding the level sensor and jamming the discharge valve. Each write is a CONTROL operation to the output area, which CARS recovers and forbids.

\begin{lstlisting}[style=carssrc,caption={The FDI overflow injection (writes to the S7 output area).},label={lst:fdi}]
c = snap7.client.Client(); c.connect(HOST, 0, 1)   # HOST = PLC 192.168.2.10
zero = bytearray(struct.pack('>f', 0.0))
while time.time() < end:
    c.write_area(PE, 0, 100, zero)   # blind LevelIn %ID100 = 0 (sensor OB30 trusts)
    c.write_area(PA, 0, 104, zero)   # jam DischargeValve %QD104 = 0 (output area 0x82)
\end{lstlisting}

\subsection*{Evaluation and proof harnesses}

\noindent \textbf{Conformance harness (\href{\repourl/blob/main/06_Build/cars_eval.py}{\texttt{cars\_eval.py}}).} Drives the labelled case set of Section~\ref{sec:evalaccuracy} through the deployed engine over \texttt{/cars/respond}, scores each answer against its intent label, and writes \texttt{cars\_eval\_matrix.csv}. It disarms for the decision-only pass so nothing is enforced on the live conduits, then re-arms.

\begin{lstlisting}[style=carssrc,caption={Driving each case through the engine and scoring it against intent.},label={lst:eval}]
def respond(src,dst,op,rate=0,dpid=1):            # one case -> the deployed decision path
    return post("/cars/respond",{"src":src,"dst":dst,"op":op,"proto":"S7","dpid":dpid,"rate":rate})
def verdict(label,resp):
    permit = resp in {"ALLOW","MONITOR","REFUSE"}
    if label in ("legit","flood_exempt"): return "TN" if permit else "FP"
    if label == "attack":                 return "TP" if not permit else "FN"
    return "grey"                                 # grey rows shown, not scored
post("/cars/defense",{"on":False},auth=True)      # decision-only pass; re-armed at the end
\end{lstlisting}

\noindent \textbf{Criticality sweep (\href{\repourl/blob/main/06_Build/cars_criticality_proof.sh}{\texttt{cars\_criticality\_proof.sh}}).} Holds the actor and operation fixed and varies only the target's tier. Part~1 runs disarmed to show the bounded elevation as a pure decision; Part~2 runs armed to measure the tier-scaled quarantine ($30 + 15w$ seconds), healing the reactive rule after each fire.

\begin{lstlisting}[style=carssrc,caption={Fire a controlled case, and heal the tier-scaled reactive rule between fires.},label={lst:critproof}]
fire(){ curl -s -XPOST $API/respond -H 'Content-Type: application/json' \
        -d "{\"src\":\"$2\",\"dst\":\"$3\",\"proto\":\"$4\",\"op\":\"$5\",\"dpid\":3}"; }
heal(){ for sw in ovsgw ovs1; do sudo ovs-ofctl -O OpenFlow13 --strict del-flows "$sw" "table=1,priority=110"; done; }
arm false     # PART 1 - decision side: bounded elevation, nothing enforced
arm true      # PART 2 - response side: real quarantine, hard_timeout = 30 + cw*15 s
\end{lstlisting}

\noindent \textbf{Campaign library (\href{\repourl/blob/main/06_Build/cars_campaign_lib.sh}{\texttt{cars\_campaign\_lib.sh}}).} The shared helpers every campaign and pivot run sources: arm/disarm through the authenticated API, a cross-layer \texttt{snap} of the live state, and \texttt{capstart} which opens the three synchronised wire captures.

\begin{lstlisting}[style=carssrc,caption={Arm/disarm and the three-point capture used across the campaign.},label={lst:camplib}]
arm(){    curl -s -X POST $API/cars/defense -H "X-CARS-Token: $TOKEN" -d '{"on":true}';  }
disarm(){ curl -s -X POST $API/cars/defense -H "X-CARS-Token: $TOKEN" -d '{"on":false}'; }
capstart(){ sudo tcpdump -i "$PLCIF" -w "$d/plc1_wire.pcap" $filt &     # what reaches the PLC
            sudo tcpdump -i any     -w "$d/of_control.pcap" 'tcp port 6653' &  # the OpenFlow flow_mods
            sudo tcpdump -i snort0  -w "$d/dpi_mirror.pcap" & }          # what the detector sees
\end{lstlisting}

\noindent \textbf{Wire campaign (\href{\repourl/blob/main/06_Build/cars_wire_campaign.sh}{\texttt{cars\_wire\_campaign.sh}}, \href{\repourl/blob/main/06_Build/cars_wire_campaign_disarmed.sh}{\texttt{cars\_wire\_campaign\_disarmed.sh}}).} Runs each attack vector against the armed system (and its disarmed twin), capturing the three wire points and timestamp-aligned state from every layer. The first vector is the control-plane disarm attempt.

\begin{lstlisting}[style=carssrc,caption={The three capture points and the control-plane vector.},label={lst:wirecamp}]
sudo timeout 240 tcpdump -i "$PLCIF" -w "$D/plc1_wire.pcap" 'host 192.168.2.10 and tcp port 102' &
sudo timeout 240 tcpdump -i any     -w "$D/of_control.pcap" 'tcp port 6653' &
sudo timeout 240 tcpdump -i snort0  -w "$D/dpi_mirror.pcap" &
# V1 CONTROL-PLANE: unauth POST /cars/defense {on:false}  -> expect HTTP 401 + DENIED audit
\end{lstlisting}

\noindent \textbf{Deep end-to-end trial (\href{\repourl/blob/main/06_Build/cars_e2e.sh}{\texttt{cars\_e2e.sh}}).} Drives the full response spectrum through the real autonomous chain and records, per scenario, the four independent proofs behind Table~\ref{tab:crosslayer}: the frame on the mirror, the Snort alert, the controller decision, and the datapath enforcement.

\begin{lstlisting}[style=carssrc,caption={The four independent ground-truth layers per scenario.},label={lst:e2e}]
wirecount(){ tshark -r "$OUT/pcap/e2e_full.pcap" -Y "$1" | wc -l; }        # WIRE   : frames on the mirror
pkts(){ ovs-ofctl -O OpenFlow13 dump-flows "$1" table=1 | grep -m1 "$2" \
        | grep -oP 'n_packets=\K[0-9]+'; }                                  # OVS    : enforcement counter moved
# SNORT : delta on /var/log/snort/alert ;  CARS : tier+response in /cars/audit
\end{lstlisting}

\noindent \textbf{Packet-drop proof (\href{\repourl/blob/main/06_Build/cars_packet_proof.sh}{\texttt{cars\_packet\_proof.sh}}).} Captures the attacker's S7 frames at two synchronised points, the Snort mirror and the PLC's physical port, armed then disarmed, giving the drop-on-the-wire counts of Table~\ref{tab:wiredrop}.

\begin{lstlisting}[style=carssrc,caption={Two synchronised capture points: what the detector sees versus what reaches the PLC.},label={lst:pktproof}]
capture(){   # tag seconds
  sudo timeout "$2" tcpdump -i snort0          "host $PLC1 and tcp port 102" -w "$D/$1_mirror.pcap" &  # DPI sees
  sudo timeout "$2" tcpdump -i enx9c69d331d874 "host $PLC1 and tcp port 102" -w "$D/$1_plc1.pcap"  & } # reaches PLC
arm true; heal; capture armed 15      # then: arm false; capture disarmed 15
\end{lstlisting}

\noindent \textbf{ICS operation battery (\href{\repourl/blob/main/06_Build/cars_ics_battery.sh}{\texttt{cars\_ics\_battery.sh}}).} Fires the full ICS operation spectrum from one trusted operator (\texttt{.31}) at the Modbus unit (\texttt{.20}) and shows the three verdicts the one conduit draws, backing Table~\ref{tab:dpi}.

\begin{lstlisting}[style=carssrc,caption={Fire each operation on the same conduit and await the new decision.},label={lst:battery}]
probe(){ restore; b=$(alast); ip netns exec opns "$@" >/dev/null 2>&1; wait_dec "$b"; }
# READ -> OPERATIONAL/ALLOW | WRITE -> SENSITIVE/THROTTLE | CONTROL,DIAG,PROGRAM,ILLEGAL -> FORBIDDEN/BLOCK
\end{lstlisting}

\subsection*{Attack clients and the process model}

\noindent \textbf{S7 attack client (\href{\repourl/blob/main/06_Build/s7_write.py}{\texttt{s7\_write.py}}).} The S7 tool used across the evaluation. Its output-area write is an actuation CARS classifies as CONTROL; the \texttt{-{}-dbspoof} mode pins a false level (the sensor-spoof deception), and further modes cover monitoring, volumetric read floods and a PLC stop.

\begin{lstlisting}[style=carssrc,caption={The output-area write (actuation) and the sensor-spoof variant.},label={lst:s7write}]
def w(v): c.write_area(PA, 0, 0, bytearray([v & 0xFF]))    # write Q0.x output image (area 0x82) = actuation
# -{}-dbspoof: pin a false level, deceiving both the control loop and the HMI:
payload = struct.pack('>f', a.spoofval)                    # S7 REAL = big-endian 32-bit float
c.db_write(a.db, a.offset, payload)                        # write-var 0x05 -> CARS classifies CONTROL
\end{lstlisting}

\noindent \textbf{Modbus attack client (\href{\repourl/blob/main/06_Build/mb_attack.py}{\texttt{mb\_attack.py}}).} Crafts arbitrary Modbus/TCP function codes, including dangerous and illegal ones, to prove CARS decides on the operation and not the five-tuple. The function code sits at offset seven, exactly where the Snort DPI anchors.

\begin{lstlisting}[style=carssrc,caption={Crafting an arbitrary Modbus function code (FC at offset 7).},label={lst:mbattack}]
def frame(fc, data=b'', tid=1, unit=1):
    pdu  = bytes([fc & 0xFF]) + data
    mbap = struct.pack('>HHHB', tid, 0x0000, len(pdu)+1, unit)   # proto-id 0x0000 @off2 ; FC @off7
    return mbap + pdu
ATTACKS = {'coil':0x05, 'diag':0x08, 'program':0x2B, 'illegal':0x64}  # CONTROL / DIAG / PROGRAM / ILLEGAL
\end{lstlisting}

\noindent \textbf{Process model (\href{\repourl/blob/main/06_Build/cars_process.py}{\texttt{cars\_process.py}}).} The soft co-simulation of the water-tank bang-bang loop, driving the real Q0.3 relay on PLC1 over S7 from the authorised controller identity. It reasserts the actuator every cycle and reports readback-versus-intent mismatches, which are the moments an attacker's write momentarily reached the relay.

\begin{lstlisting}[style=carssrc,caption={The bang-bang control law, the dynamics, and the interference health metric.},label={lst:process}]
if   level <= LOW:  pump = True                 # control law (eq.1)
elif level >= HIGH: pump = False
level += FILL if pump else -DRAIN               # dynamics (eq.2)
c.write_area(PA, 0, 0, bytearray([0x08 if pump else 0x00]))   # actuate + reassert Q0.3 every cycle
if (got != 0) != pump: bad += 1                 # readback != intent = attacker interference reached the relay
\end{lstlisting}

\subsection*{Gap mitigations: DPI hardening and the process guardian}

\noindent \textbf{DPI reassembly hardening (\texttt{cars.conf}, \texttt{cars.rules}).} The single-packet S7 rules were defeated by TCP fragmentation (Section~\ref{sec:evaldpi}). Two changes closed it: stream reassembly on the ICS ports, and re-anchoring the operation match to the S7 job header instead of a fixed packet offset.

\begin{lstlisting}[style=carssrc,caption={Stream reassembly, and the offset-to-relative rule change.},label={lst:dpihard}]
# cars.conf : reassemble the ICS conduits before inspection
preprocessor stream5_global: track_tcp yes, track_udp no, max_tcp 8192
preprocessor stream5_tcp:    policy linux, timeout 60, ports both 102 502
# cars.rules S7 Write-Var, before (fixed offset) -> after (PDU-relative):
#  before: content:"|03 00|"; depth:2; content:"|32 01|"; offset:7; depth:2; content:"|05|"; offset:17; depth:1;
#  after : content:"|03 00|"; content:"|32 01|"; distance:5; within:7; content:"|05|"; distance:8; within:9;
\end{lstlisting}

\noindent \textbf{Fragmentation harness (\href{\repourl/blob/main/report/gap_evidence/gap_evidence/scripts/frag_s7_write.py}{\texttt{frag\_s7\_write.py}}).} Opens a real allowlisted S7 session and sends a Write-Var either whole or split across TCP segments so the function byte falls in the second, the test behind the before/after of Section~\ref{sec:evaldpi}.

\begin{lstlisting}[style=carssrc,caption={Sending the Write-Var whole or split at byte N (0x05 pushed into the second segment).},label={lst:frag}]
s.setsockopt(socket.IPPROTO_TCP, socket.TCP_NODELAY, 1)   # Nagle off: each send() is its own segment
s.sendall(COTP_CR); s.recv(512); s.sendall(S7_SETUP); s.recv(512)   # real S7 session on an allowlisted conduit
if split <= 0: s.send(WRITEVAR)                            # whole  -> detected -> ISOLATE
else:          s.send(WRITEVAR[:split]); time.sleep(gap); s.send(WRITEVAR[split:])  # fragmented
\end{lstlisting}

\noindent \textbf{Low-rate insider tools (\href{\repourl/blob/main/report/gap_evidence/gap_evidence/scripts/hmi_cmd.py}{\texttt{hmi\_cmd.py}}, \href{\repourl/blob/main/report/gap_evidence/gap_evidence/scripts/spoof_high.py}{\texttt{spoof\_high.py}}).} A single authorised HMI command, and a sustained high sensor spoof: the two residual insider actions of Section~\ref{sec:evalguardian}.

\begin{lstlisting}[style=carssrc,caption={One authorised HMI stop; and pinning the level sensor high.},label={lst:insider}]
# hmi_cmd.py: one authorised write of the command merker %MB0 (bit 2 = HMI_Stop)
c.write_area(MK, 0, 0, bytearray([0x04]))                       # CARS classifies CONTROL -> ALLOW (trusted role)
# spoof_high.py: pin %ID100 high so Sim.Level (= LevelIn*20) reads ~full
c.write_area(PE, 0, 100, bytearray(struct.pack('>f', 4.5)))    # reported ~90 while the real tank drains
\end{lstlisting}

\noindent \textbf{Process guardian (\href{\repourl/blob/main/report/gap_evidence/gap_evidence/scripts/cars_process_guardian.py}{\texttt{cars\_process\_guardian.py}}).} The defence-in-depth layer of Section~\ref{sec:evalguardian}: an independent, read-only monitor that adds the high-side envelope, symmetric rate and liveness checks the low-side remediation agent lacks, and raises alarms only. It never writes the PLC and never touches enforcement or the remediation agent.

\begin{lstlisting}[style=carssrc,caption={The high-side envelope, rate and liveness checks (alarm-only).},label={lst:guardian}]
if lvl > CEIL:    alarm("HIGH_LEVEL", level=lvl, ceiling=CEIL)   # the excursion the low-side agent MISSES
if lvl < FLOORV:  alarm("LOW_LEVEL",  level=lvl, floor=FLOORV)
if prev is not None and abs(lvl - prev) > MAXDELTA:
    alarm("RATE", level=lvl, prev=prev)                          # a physically impossible per-poll jump
if abs(lvl - last_val) <= 0.6 and now - last_change > STALL_S:
    alarm("STALL", level=lvl)                                    # level frozen while it should cycle = stopped/held
\end{lstlisting}

\section{The CARS intelligence dashboard}
\label{appx:dashboard}

The dashboard (\href{\repourl/blob/main/06_Build/cars_dashboard.py}{\texttt{cars\_dashboard.py}}) is a live operator view of the defence rather than a static diagram. It is discovery-driven: a few times a second it polls the controller's REST interface (\texttt{/cars/status}, \texttt{/cars/hosts}, \texttt{/cars/ports}, \texttt{/cars/links}, \texttt{/cars/guard}, \texttt{/cars/audit}) and renders whatever the controller currently knows. The nodes (switches, hosts, the controller, and the modelled north side of IT, DMZ, firewall and Snort), together with their port bindings, link health, the anti-spoof GUARD drop counters and the engine's own decisions, are all pulled from the running system, so the picture is the live state and not a description of it. A separate pane reads the remediation agent's status and event feed from its local files, because that agent runs inside the OT namespace and cannot reach the control plane.

\begin{figure}[t]
\centering
\includegraphics[width=0.9\textwidth]{dash_overview.png}
\caption{The CARS dashboard's discovered topology. Switches (\texttt{ovsgw}, \texttt{ovs1}, \texttt{ovs2}) in blue, the supervisory hosts (EWS, SCADA, the historians, CARS-Remediation) and PLCs discovered live with their roles and criticality rings, the CARS controller and Snort IDS, the modelled north side (Enterprise firewall, DMZ/jump, OT NAT) and the Cell-2 NAT path to PLC2. Unregistered hosts (the on-segment attacker seams \texttt{192.168.2.66} and \texttt{192.168.2.67}) are flagged. Everything shown is polled live from the controller, not drawn by hand.}
\label{fig:dashtopo}
\end{figure}

The centre of the view is the discovered network topology (Figure~\ref{fig:dashtopo}); the right-hand column carries the live decision feed, filterable by criticality tier, response, operation class and enforcement mode. Each row shows the source and its discovered role, the destination and its criticality tier, the operation class, the criticality-elevated tier and the selected response, colour-coded by severity. It is this feed that makes the engine's reasoning legible in real time: a supervisory read appears as an OPERATIONAL ALLOW, a remediation write as a RESTORE, and an out-of-role actuation as a FORBIDDEN ISOLATE.

Figure~\ref{fig:dashisolate} in Section~\ref{sec:evalcampaign} is a capture of this feed during the armed false-data-injection run: the compromised supervisory host \texttt{192.168.2.31 (scada)} issuing S7 CONTROL to the CRITICAL PLC1, each attempt classed FORBIDDEN and met with an ENFORCED ISOLATE. Placed beside the datapath rule of Figure~\ref{fig:p6armedflow} and the wire capture of Section~\ref{sec:evalcrosslayer}, it is the same attack seen at a third layer, the operator's, completing the cross-layer account: one event visible at once as a policy decision, an installed flow and a dropped packet.

\section{Gap remediation evidence}
\label{appx:gapevidence}

This appendix collects the fresh before-and-after evidence for the late-stage gap remediations of Chapter~\ref{chap:evaluation}, captured on 7 and 8 August 2026 on the armed testbed and returned to green after each run. The consolidated record with reproduction commands is the companion file \texttt{GAP\_REMEDIATION\_EVIDENCE.md}; the raw artefacts (controller-audit slices, the \texttt{cars.rules}/\texttt{cars.conf} diffs, the guardian log and the honeypot captures) are in \texttt{gap\_evidence/}.

\noindent\textbf{Gap A: DPI defeated by fragmentation, then hardened.} A Write-Var split across two TCP segments evaded the fixed-offset single-packet rule; after enabling stream reassembly and re-anchoring the rules to the S7 job header it is recovered and isolated at more than one split point, with the legitimate loop still permitted.
\begin{lstlisting}[style=carsdump,caption={Gap A: before (fragmented write evades) and after (caught), from the controller audit.},label={lst:ev_dpi}]
BEFORE (single-packet rules) - fragmented at byte 14, armed:
  192.168.2.31(scada) -> 192.168.2.10(plc) TCP => ALLOW (operational)   [no CONTROL recovered; 0xca=0; write reached PLC]
  (same write sent WHOLE):
  FORBIDDEN 192.168.2.31(scada) -> 192.168.2.10(plc) S7 CONTROL => ISOLATE source 192.168.2.31 75s
AFTER (stream5 reassembly + PDU-relative rules) - same fragmented attack:
  --split 14: FORBIDDEN 192.168.2.31 ... S7 CONTROL => ISOLATE source 192.168.2.31 75s ; flow 0xca nw_src=192.168.2.31
  --split  9: FORBIDDEN 192.168.2.31 ... S7 CONTROL => ISOLATE
  legit loop: OPERATIONAL 192.168.2.55(ews) -> 192.168.2.10 READ/CONTROL => ALLOW   [no new false positive]
\end{lstlisting}

The fix itself is shown as the unified diff of the deployed Snort configuration and rules against the pre-hardening backup: Listing~\ref{lst:ev_confdiff} adds the TCP stream-reassembly preprocessor on the ICS ports, and Listing~\ref{lst:ev_rulesdiff} re-anchors every S7 operation rule from a fixed packet offset to a PDU-relative match on the S7 job header.
\lstinputlisting[style=carsdump,caption={Gap A: \texttt{cars.conf} diff, stream reassembly added on ports 102 and 502.},label={lst:ev_confdiff}]{appendix_listings/gapA_cars_conf.diff}
\lstinputlisting[style=carsdump,caption={Gap A: \texttt{cars.rules} diff, S7 rules re-anchored PDU-relative (rev 1 to rev 2).},label={lst:ev_rulesdiff}]{appendix_listings/gapA_cars_rules.diff}

\noindent\textbf{Gap B: trusted-insider process guardian.} The high-rate overflow is cut by the volumetric overlay even from a trusted role; the read-only guardian then catches the low-rate denial and the high-side sensor spoof that the low-side remediation agent misses.
\begin{lstlisting}[style=carsdump,caption={Gap B: the flood backstop and the two guardian alarms.},label={lst:ev_guardian}]
FLOOD BACKSTOP (trusted .45 driven at high rate):
  192.168.2.45(remediation) -> 192.168.2.10 S7 CONTROL => [FLOOD 51 ops/s] BLOCK conduit 75s  [FDI cut after 77 writes; tank held ~mid-70s]
DEMO 1 - low-rate denial (one authorised HMI_Stop from .45):
  audit:    192.168.2.45(remediation) -> 192.168.2.10 S7 CONTROL => ALLOW (operational)   [network correctly permits]
  guardian: [GUARDIAN] ALARM STALL {level: 40.4, frozen_s: 12.0}                          [process halt caught]
DEMO 2 - spoof-to-high (reported level pinned at 90, real tank drains):
  remediation status: level=90.0, restores unchanged                                      [low-side agent silent = blind]
  guardian: [GUARDIAN] ALARM HIGH_LEVEL {level: 90.0, ceiling: 78.0}  +  ALARM RATE        [high-side caught]
\end{lstlisting}

The guardian writes each alarm to a read-only feed; the stall record from the fresh run is the raw line in Listing~\ref{lst:ev_guardianfeed}. This is a crafted scenario that demonstrates the mechanism rather than broad operational resilience: the flood, envelope and rate thresholds are tuned to this tank and its loop rate, and would need re-tuning for a different plant, loop cadence or operator profile.
\lstinputlisting[style=carsdump,caption={Gap B: the guardian alarm feed (\texttt{cars\_guardian.jsonl}), stall record from the fresh run.},label={lst:ev_guardianfeed}]{appendix_listings/gapB_guardian.jsonl}

\noindent\textbf{Gap C: flow-integrity poll window.} A rule injected and deleted inside the ten-second poll evades the auditor; a persistent one is caught within one poll. The single traces below were then repeated thirty times each to quantify the window: the persistent case was caught on every trial at a median of 8.5\,seconds, the two-second transient on only five of thirty (the runs whose life straddled a poll boundary). The per-trial values are in \texttt{gap\_evidence/flowint/persistent.csv} and \texttt{transient.csv}.
\begin{lstlisting}[style=carsdump,caption={Gap C: transient (mostly missed) versus persistent (always caught) flow-integrity status, with the thirty-trial totals.},label={lst:ev_poll}]
baseline:                     {"ok": 1, "extra": 0}
2 s transient inject+delete:  t+7s ok:1 | t+14s ok:1 | t+21s ok:1     [MISSED - never flipped]
persistent inject (left):     {"ok": 0, "extra": 1}                   [CAUGHT within one poll]
after cleanup:                {"ok": 1, "extra": 0}
30-trial totals:  persistent 30/30 caught, median 8.5 s | transient 5/30 caught, 25/30 evaded
\end{lstlisting}

The event-driven flow-monitor of Section~\ref{sec:threats} was then run against the identical thirty-trial transient test. It subscribes to the switch's Nicira flow-monitor (\texttt{ovs-ofctl monitor <br> watch:}) and, on any change to the policy tables that is not a cookie-\texttt{0xca} reactive rule, re-runs the deployed auditor's trusted check at once, so the sub-poll window closes without changing the verdict logic. Every one of the thirty two-second injections was caught, at a median of 0.27\,seconds, against five of thirty for the ten-second poll; the per-trial values are in \texttt{gap\_evidence/flowint/transient\_eventdriven.csv}.
\begin{lstlisting}[style=carsdump,caption={Gap C remedy: event-driven monitor on the same transient test.},label={lst:ev_poll_fix}]
10 s poll (baseline):     transient  5/30 caught, 25/30 evaded (median 8.5 s when caught)
event-driven monitor:     transient 30/30 caught, median 0.27 s   [sub-poll gap closed]
mechanism: NXST_FLOW_MONITOR watch -> immediate trusted re-check on any non-0xca policy-table change
\end{lstlisting}

\noindent\textbf{Gap 2: conduit block at a shared / NAT-collapsed identity.} An IT attacker source-NATed onto the gateway identity \texttt{192.168.2.1} was originally cut by a source isolation, which quarantines every host behind the NAT (Section~\ref{sec:evalanatomy}). The response selection was refined so a shared identity is cut at the conduit instead. Verified on the deployed engine: the same forbidden attempt from \texttt{192.168.2.1} now returns \texttt{BLOCK} and installs a conduit drop, while a true single host is still source-isolated.
\begin{lstlisting}[style=carsdump,caption={Gap 2: shared-identity conduit BLOCK versus single-host source ISOLATE, deployed engine.},label={lst:ev_natblock}]
gateway .1 (NAT-collapsed) -> CONTROL -> PLC1:  response=BLOCK
  flow: cookie=0xca priority=100,ip,nw_src=192.168.2.1,nw_dst=192.168.2.10 actions=drop   [conduit: only .1->PLC]
unknown .77 (true host)   -> CONTROL -> PLC1:  response=ISOLATE
  flow: cookie=0xca priority=110,ip,nw_src=192.168.2.77 actions=drop                       [source: whole host]
both land zero writes on the PLC; the gateway's non-PLC traffic is untouched by the conduit match.
\end{lstlisting}

\noindent\textbf{Gap D: DEFLECT diverts to the honeypot.} Forcing DEFLECT installs the redirect and the attacker's PLC-bound traffic arrives at the decoy while the real PLC sees nothing.
\begin{lstlisting}[style=carsdump,caption={Gap D: silent-drop baseline, the redirect, and the decoy receiving the diverted traffic.},label={lst:ev_deflect}]
silent-drop baseline:  atkns(.66) ping 192.168.2.10 -> 100% loss                       [dropped; attacker hangs]
force DEFLECT:         response=DEFLECT  action="DEFLECT conduit -> honeypot 192.168.3.99 (deception, self-healing)"
redirect flow:         priority=105 nw_src=192.168.2.66 nw_dst=192.168.2.10 -> set ip_dst=192.168.3.99, goto t2 (n_packets=3)
honeypot capture:      IP 192.168.2.66 > 192.168.3.99: ICMP echo request               [attacker DIVERTED to decoy]
                       ARP, who-has 192.168.2.66 tell 192.168.3.99                      [decoy engaging]
\end{lstlisting}

\noindent This is a proof of concept for the redirect path, not evidence of robustness. A determined adversary who notices the diversion, fingerprints the decoy or changes tactics is outside what this test exercises; sustaining a convincing interactive decoy is carried as future work in Section~\ref{sec:futurework}.

\noindent The plaintext control channel (G1), state-exhaustion resilience (G3) and the file-based detection bridge under heavy load (G4) are verified but framed as bounded limitations in Section~\ref{sec:threats} rather than tested, for the reasons given there.

\section{Process-transparency and availability runs}
\label{appx:availability}

This section holds the raw artefacts behind the process-transparency claim of Section~\ref{sec:evalprocessrisk} and the reconnection test of Section~\ref{sec:threats}, captured on the armed testbed on 7 August 2026 and returned to green after each run.

\noindent\textbf{Transparency: armed against disarmed under matched load.} The two states were interleaved in short blocks over two hours so that any drift in the hardware-in-the-loop simulation fell on both equally. The per-second tank level, phase label, online flag and restore count are in \texttt{availability/interleaved.csv} (7{,}060 samples); the online-only summary is below. The means and spreads coincide and neither state left the control band, so arming the defence has no measurable effect on the process. As Section~\ref{sec:evalprocessrisk} states, this is strong evidence for the testbed and its hardware-in-the-loop tank rather than a universal guarantee for all industrial processes.
\begin{lstlisting}[style=carsdump,caption={Transparency: online-only level statistics per phase, from \texttt{interleaved.csv}.},label={lst:ev_transparency}]
disarmed:  n=3528  mean=51.5%  sd=12.8  min=25.5  max=71.3   (band 20-78, no excursion)
armed:     n=3530  mean=50.7%  sd=13.2  min=25.6  max=71.3   (band 20-78, no excursion)
EWS/SCADA legitimate throughput unchanged across both phases.
\end{lstlisting}

\noindent\textbf{Block-and-maintain and self-heal, measured at the process.} An allowlisted supervisory host launched a false-data injection at the CRITICAL PLC under the armed defence. The write was classified FORBIDDEN and the source isolated after nine packets with the criticality-scaled seventy-five-second timeout; the tank stayed within band for the whole block, the legitimate conduits kept flowing, and the restoration agent was never called. The rule then expired and the source returned to service with no operator action, the install-to-withdrawal spanning seventy-six seconds.
\begin{lstlisting}[style=carsdump,caption={Block-and-maintain under attack, then self-heal, from the armed run.},label={lst:ev_selfheal}]
t+0.0s   FDI write from 192.168.2.31(scada) -> 192.168.2.10(plc, CRITICAL) => FORBIDDEN
t+~0s    ISOLATE source 192.168.2.31 (cookie 0xca, hard_timeout=75s, both switches)
block    tank level held within band; legit EWS/HMI conduits uninterrupted; restore agent NOT called
t+76s    rule expired (SEND_FLOW_REM); source restored, no operator action; table back to baseline
\end{lstlisting}

\noindent\textbf{Reconnection jitter, two planes.} Restarting the passive Snort sensor left the live conduit's inter-frame timing indistinguishable from baseline; forcing a controller reconnection by delete-and-re-add opened a roughly twenty-second window with multi-second forwarding stalls and dropped the synchronous S7 session. The per-second timeline (level, legitimate packets per second, controller-attached flag, event marks) is in \texttt{availability/jitter\_timeline.csv} with the accompanying \texttt{hil.pcap}. As Section~\ref{sec:threats} sets out, that window largely reflects the heavy re-add and pipeline re-sync rather than a faithful crash against the \texttt{fail\_mode=secure} switches, so it is reported as an upper bound and the clean crash-transparency figure is left as future work. This measurement should not be read as a crash-resilience proof or a production-availability guarantee: a fail-secure switch under a forced delete-and-re-add is not the same as a real controller crash, and the deployment implication is left deliberately cautious.
\begin{lstlisting}[style=carsdump,caption={Reconnection jitter: the two planes contrasted.},label={lst:ev_jitter}]
Snort restart (passive, out-of-band):  inter-frame mean 4.90 ms vs 4.91 ms baseline, worst gap 33 ms vs 42 ms  [TRANSPARENT]
Controller del/re-add (forced):         ~20 s window, worst forwarding gap 4.2 s, S7 session dropped, manual reconnect  [UPPER BOUND, not a faithful crash]
\end{lstlisting}

\section{Statistical visualisations of the evaluation}
\label{appx:stats}

The figures below plot the measured data behind the Critical Evaluation. Each is generated directly from a captured result file, named in its caption, with no smoothing or illustrative curves; log axes are used where a range spans orders of magnitude.

\begin{figure}[h]
\centering
\includegraphics[width=0.95\textwidth]{stat_traffic_baseline.png}
\caption{Normal OT traffic baseline: per-conduit mean packets/s and the protocol mix, from a 31-second steady-state capture (6{,}490 packets, 9 conduits). The cyclic engineering-to-PLC S7 flow dominates; the remainder is PLC replies and background multicast. Source: \texttt{results/b1/baseline.pcap} via \href{\repourl/blob/main/07_Evaluation/overnight/b1_analyze.py}{\texttt{b1\_analyze.py}}.}
\label{fig:stattraffic}
\end{figure}

\begin{figure}[h]
\centering
\includegraphics[width=0.62\textwidth]{stat_confusion.png}
\caption{Decision conformance at scale over the 2{,}078-case labelled corpus, which is attack-heavy by design (2{,}040 attacks against 38 legitimate or exempt cases). Every attack was restrained (no false negative) and every legitimate or exempt case permitted (no false positive). Because the decision is deterministic and the corpus is imbalanced, the two per-class error rates, both zero, are the meaningful figures rather than any aggregate. The attack side, where the label and the default-deny share the same registration test, is a coverage-and-robustness result; the zero false positive over the authorised cases is the trust result. Source: \texttt{results/vast/vast.csv}.}
\label{fig:statconfusion}
\end{figure}

\begin{figure}[h]
\centering
\includegraphics[width=0.98\textwidth]{stat_criticality_surface.png}
\caption{Operation- and source-graded response over 240 live judgements. The trusted host keeps its reads (with a small tail of persistence-escalated cuts as its offence count rises across the sequence) and is cut on the operations its role must not perform, whereas the unregistered attacker is cut on every operation. Source: \texttt{results/critbeh/judgements.csv}.}
\label{fig:statcrit}
\end{figure}

\begin{figure}[h]
\centering
\includegraphics[width=0.95\textwidth]{stat_flowintegrity.png}
\caption{Flow-integrity, ten-second poll versus the event-driven monitor. Left: a two-second transient injection is missed by the poll on 25 of 30 trials but caught on all 30 by the event-driven monitor; a persistent injection is caught by both. Right: median detection latency, 8.5\,s for the poll against 0.27\,s for the monitor (log axis). Sources: \texttt{results/flowint/\{persistent,transient\}.csv}, \texttt{results/gap4\_flowmonitor/transient\_eventdriven.csv}.}
\label{fig:statflowint}
\end{figure}

\begin{figure}[h]
\centering
\includegraphics[width=0.8\textwidth]{stat_transparency.png}
\caption{Process transparency: the tank-level distribution armed against disarmed, over roughly seven thousand interleaved, load-matched per-second samples. The two distributions coincide (means 51.5 against 50.7\%) and neither leaves the control band, so arming the defence has no measurable effect on the process. Source: \texttt{results/e2/interleaved.csv}.}
\label{fig:stattransparency}
\end{figure}

\begin{figure}[h]
\centering
\includegraphics[width=0.78\textwidth]{stat_mttm_load.png}
\caption{Reaction window against background load, measured at the controller's decision-and-log stage: both the background and the probe are posted directly to \texttt{/cars/respond}, so this characterises the controller (\texttt{classify}, \texttt{select\_response} and the audit write), not the Snort-to-file-to-HTTP sensor bridge. The window holds in the tens of milliseconds to about a thousand operations a second and rises only as this single-threaded stage nears its sixteen-to-eighteen-thousand-a-second bench ceiling. The sensor bridge is a separate and lower ceiling, the acknowledged bottleneck carried to the future work (Section~\ref{sec:threats}), and is not what this curve measures. Source: \texttt{results/gap1\_live/curve.csv}.}
\label{fig:statload}
\end{figure}

\begin{figure}[h]
\centering
\includegraphics[width=0.92\textwidth]{stat_jitter.png}
\caption{Reconnection jitter on the legitimate conduit: restarting the passive Snort sensor is transparent (no dip), whereas a forced controller reconnection stalls forwarding. Source: \texttt{results/jitter/timeline.csv} (one spurious counter sample removed).}
\label{fig:statjitter}
\end{figure}

\subsection{Long-run and stress-test campaign}
\label{appx:longrun}

The figures below back the long-run and stress test of Section~\ref{sec:evallongrun}: four autonomous runs on the one testbed, 440 measured attacks, all enforced. They extend the earlier evidence in duration and repetition, not in network size. The four runs are repeated runs on the same testbed and machine state, not independent replicates; where a median is quoted it is pooled across the runs rather than averaged across per-run medians. Each plot is generated directly from the per-run \texttt{logs/*.csv} by \href{\repourl/blob/main/07_Evaluation/overnight/night/night_analyze.py}{\texttt{night\_analyze.py}}, with the row counts and any exclusions stated in the caption.

\begin{figure}[h]
\centering
\includegraphics[width=0.95\textwidth]{stat_longrun_mttm.png}
\caption{Reaction window over the long-run and stress campaign. The source-cut path (356 attacks, median 8.5\,ms) sits close to the 7.6\,ms controlled single-clock reference of Section~\ref{sec:evallatency}; the false-data injections from the compromised trusted host form a separate mode near 100\,ms, the conduit-block path that installs after the first packets. 437 of the 440 measured attacks are shown: two under-load DDoS-probe samples, which belong to the load curve rather than the clean reaction window, and one physically impossible negative row from a stale-rule read, are excluded. Source: the four \texttt{night\_*/logs/mttm\_all.csv}.}
\label{fig:statlongmttm}
\end{figure}

\begin{figure}[h]
\centering
\includegraphics[width=0.95\textwidth]{stat_longrun_pervector.png}
\caption{Per-vector reaction window pooled over the four runs (box is the interquartile range, line the median). The source-cut vectors sit at roughly 7 to 12\,ms; the two false-data injection vectors from the compromised trusted host sit near 100\,ms because they draw a conduit block rather than an immediate source cut. Source: \texttt{night\_*/logs/mttm\_all.csv}.}
\label{fig:statlongpervec}
\end{figure}

\begin{figure}[h]
\centering
\includegraphics[width=0.9\textwidth]{stat_longrun_tierladder.png}
\caption{Criticality-scaled isolation timeout over 120 tier probes across the four runs. Every probe installed 75/60/45/30\,s for CRITICAL/HIGH/MEDIUM/LOW with no deviation, at a median engine decide-time of 0.50\,ms. Source: \texttt{night\_*/logs/tierladder.csv}.}
\label{fig:statlongtier}
\end{figure}

\begin{figure}[h]
\centering
\includegraphics[width=0.95\textwidth]{stat_connpool_gradient.png}
\caption{The S7-1200 (1212C) connection-pool limit. The CPU exposes six dynamic connection resources; repeatedly isolating S7 sessions leaves half-open sessions that consume them. Packing about five isolations into each short cycle evicted the legitimate Factory IO client at 4.0 times an hour and the run could not be sustained; pacing to one isolation per cycle held the process for about ninety minutes at 1.3 to 2.0 evictions an hour. This is the empirical confirmation of the connection-pool caveat in Section~\ref{sec:threats}. Source: \texttt{night\_*/logs/\{freeze,monitor,mttm\_all\}.csv}.}
\label{fig:statconnpool}
\end{figure}

That the freeze is the endpoint's limit rather than a defect in the policy is confirmed on the switch: an isolate installs a single drop scoped to the attacker's source address alone, so it cannot match the legitimate conduit. Listing~\ref{lst:connscope} is a live capture taken on \texttt{ovs1} while an isolation was in force, showing the reactive rule matching only \texttt{nw\_src=192.168.2.77} while the \texttt{192.168.2.55} to \texttt{192.168.2.10} conduit remained a committed allow with its counter advancing.

\lstinputlisting[style=carsdump,caption={The isolate is source-scoped: the reactive drop matches the attacker alone, and the legitimate Factory IO conduit is untouched during the isolation (captured live on \texttt{ovs1}).},label={lst:connscope}]{appendix_listings/conn_scope.txt}

\begin{figure}[h]
\centering
\includegraphics[width=0.98\textwidth]{stat_longrun_stability.png}
\caption{Process stability under continuous attack: the tank level across the ninety-two-minute paced run, held within the 20 to 78 control band throughout. The controller and its remediation agent reported the process online in all ninety-two one-minute snapshots (this is the agent's PLC link sampled once a minute, which does not capture sub-minute client outages) and the flow auditor reported no drift. The two dashed marks are the brief Factory IO client evictions that the liveness guard caught and recovered. Source: \texttt{night\_paw2\_0153/logs/\{monitor,freeze\}.csv}.}
\label{fig:statlongstab}
\end{figure}
